\textbf{Phần 1. Trắc nghiệm lựa chọn}
\setcounterpageref{ex}{0}

\begin{ex} %[4D7V33]
Số \textit{Hai mươi ba triệu bốn trăm linh năm nghìn} viết là
\choice
{$23\,450\,000$}
{\True $23\,405\,000$}
{$23\,045\,000$}
{$234\,005\,000$}
\loigiai{
Số đó là $23\,405\,000$.\\
\textbf{Đáp án của câu hỏi là B}
}
\end{ex}

\begin{ex} %[4D7V33]
Chữ số 5 trong số $75\,200\,140$ thuộc hàng nào?
\choice
{\True Hàng triệu}
{Hàng trăm nghìn}
{Hàng chục triệu}
{Hàng chục nghìn}
\loigiai{
Chữ số 5 ở hàng triệu.\\
\textbf{Đáp án của câu hỏi là A}
}
\end{ex}

\begin{ex} %[4D7V33]
Làm tròn số $4\,680\,000$ đến hàng trăm nghìn ta được
\choice
{$4\,600\,000$}
{$5\,000\,000$}
{\True $4\,700\,000$}
{$4\,680\,000$}
\loigiai{
Hàng chục nghìn là 8 nên làm tròn thành $4\,700\,000$.\\
\textbf{Đáp án của câu hỏi là C}
}
\end{ex}

\begin{ex} %[4D7V33]
Số lớn nhất có 7 chữ số khác nhau là
\choice
{$9\,999\,999$}
{$9\,876\,540$}
{$9\,867\,543$}
{\True $9\,876\,543$}
\loigiai{
Số lớn nhất có 7 chữ số khác nhau là $9\,876\,543$.\\
\textbf{Đáp án của câu hỏi là D}
}
\end{ex}

\begin{ex} %[4D7V33]
Số liền sau của số $19\,999\,999$ là
\choice
{$19\,999\,998$}
{\True $20\,000\,000$}
{$20\,000\,001$}
{$21\,000\,000$}
\loigiai{
$19\,999\,999 + 1 = 20\,000\,000$.\\
\textbf{Đáp án của câu hỏi là B}
}
\end{ex}

\begin{ex} %[4D7V33]
Trong số $345\,678\,900$, các chữ số thuộc lớp nghìn là
\choice
{\True $6; 7; 8$}
{$3; 4; 5$}
{$9; 0; 0$}
{$3; 6; 9$}
\loigiai{
Lớp nghìn gồm: 6, 7, 8.\\
\textbf{Đáp án của câu hỏi là A}
}
\end{ex}

\begin{ex} %[4D7V33]
Sắp xếp các số sau theo thứ tự từ bé đến lớn: $35\,020\,000; 35\,200\,000; 35\,002\,000; 35\,220\,000$.
\choice
{$35\,200\,000 < 35\,020\,000 < 35\,002\,000 < 35\,220\,000$}
{$35\,002\,000 < 35\,200\,000 < 35\,020\,000 < 35\,220\,000$}
{\True $35\,002\,000 < 35\,020\,000 < 35\,200\,000 < 35\,220\,000$}
{$35\,220\,000 < 35\,200\,000 < 35\,020\,000 < 35\,002\,000$}
\loigiai{
Thứ tự: $35\,002\,000 < 35\,020\,000 < 35\,200\,000 < 35\,220\,000$.\\
\textbf{Đáp án của câu hỏi là C}
}
\end{ex}

\begin{ex} %[4D7V33]
Tìm chữ số $x$ để $45,x78,000 > 45\,878\,000$:
\choice
{$x = 8$}
{$x = 7$}
{$x = 6$}
{\True $x = 9$}
\loigiai{
$x = 9$ thỏa mãn $45\,978\,000 > 45\,878\,000$.\\
\textbf{Đáp án của câu hỏi là D}
}
\end{ex}

\begin{ex} %[4D7V33]
Viết số $M = 70\,000\,000 + 400\,000 + 5000 + 60$:
\choice
{$74\,005\,060$}
{\True $70\,405\,060$}
{$70\,450\,060$}
{$740\,005\,060$}
\loigiai{
$70\,405\,060$.\\
\textbf{Đáp án của câu hỏi là B}
}
\end{ex}

\begin{ex} %[4D7V33]
Số nào sau đây chia hết cho cả 2 và 5?
\choice
{\True $24\,680$}
{$12\,455$}
{$35\,791$}
{$46\,802$}
\loigiai{
Số tận cùng bằng 0 chia hết cho cả 2 và 5.\\
\textbf{Đáp án của câu hỏi là A}
}
\end{ex}

\begin{ex} %[4D7V33]
Số bé nhất có 8 chữ số khác nhau có chữ số hàng chục triệu là 1 là
\choice
{$10\,000\,000$}
{$12\,345\,678$}
{\True $10\,234\,567$}
{$10\,234\,576$}
\loigiai{
$10\,234\,567$.\\
\textbf{Đáp án của câu hỏi là C}
}
\end{ex}

\begin{ex} %[4D7V33]
Có bao nhiêu số có 7 chữ số mà tổng các chữ số của nó bằng 2?
\choice
{$6$ số}
{$8$ số}
{$9$ số}
{\True $7$ số}
\loigiai{
Các số đó là: $2\,000\,000$ (1 số) và các số bắt đầu bằng 1 và có thêm 1 chữ số 1 ở 6 vị trí còn lại (6 số). Tổng cộng có $1 + 6 = 7$ số.\\
\textbf{Đáp án của câu hỏi là D}
}
\end{ex}

\textbf{Phần 2. Trắc nghiệm đúng sai}
\setcounterpageref{ex}{0}

\begin{ex} %[4D7V33]
Xét tính đúng sai về hệ thống hàng và lớp:
\choiceTF
{\True Lớp triệu gồm hàng triệu, hàng chục triệu, hàng trăm triệu.}
{\True Số $100\,000\,000$ là số bé nhất có 9 chữ số.}
{Số $45\,678\,900$ có chữ số 5 ở hàng trăm nghìn.}
{\True Số $45\,678\,900$ có chữ số 5 ở hàng triệu.}
\loigiai{
\begin{itemize}
    \item Lớp triệu: Đúng.
    \item Số có 9 chữ số bé nhất: Đúng.
    \item Chữ số 5 ở hàng triệu (không phải trăm nghìn). Ý 3 Sai, Ý 4 Đúng.
\end{itemize}
\begin{itemize}
    \item a) \textbf{Đúng}.
    \item b) \textbf{Đúng}.
    \item c) \textbf{Sai}.
    \item d) \textbf{Đúng}.
\end{itemize}
}
\end{ex}

\begin{ex} %[4D7V33]
Cho số $864\,200\,000$. Xét tính đúng sai:
\choiceTF
{\True Số này đọc là Tám trăm sáu mươi tư triệu hai trăm nghìn.}
{\True Lớp triệu là $864$.}
{Chữ số 4 có giá trị là $400\,000$.}
{\True Chữ số 4 có giá trị là $4\,000\,000$.}
\loigiai{
\begin{itemize}
    \item Đọc số đúng.
    \item Lớp triệu: 864. Đúng.
    \item Chữ số 4 ở hàng triệu nên có giá trị $4\,000\,000 \neq 400\,000$. Ý 3 Sai, Ý 4 Đúng.
\end{itemize}
\begin{itemize}
    \item a) \textbf{Đúng}.
    \item b) \textbf{Đúng}.
    \item c) \textbf{Sai}.
    \item d) \textbf{Đúng}.
\end{itemize}
}
\end{ex}

\textbf{Phần 3. Trả lời ngắn}
\setcounterpageref{ex}{0}

\begin{ex} %[4D7V33]
Viết số: Bốn mươi lăm triệu sáu trăm nghìn.
\shortans{45600000}
\loigiai{
\textbf{Đáp số:} $45600000$\\
$45\,600\,000$.
}
\end{ex}

\begin{ex} %[4D7V33]
Làm tròn số $6\,350\,000$ đến hàng trăm nghìn.
\shortans{6400000}
\loigiai{
\textbf{Đáp số:} $6400000$\\
$6\,400\,000$.
}
\end{ex}

\begin{ex} %[4D7V33]
Trong số $567\,890\,123$, tính tổng các chữ số thuộc lớp triệu.
\shortans{18}
\loigiai{
\textbf{Đáp số:} $18$\\
$5 + 6 + 7 = 18$.
}
\end{ex}

\begin{ex} %[4D7V33]
Tìm số liền trước của số $30\,000\,000$.
\shortans{29999999}
\loigiai{
\textbf{Đáp số:} $29999999$\\
$29\,999\,999$.
}
\end{ex}

\begin{ex} %[4D7V33]
Có bao nhiêu số chẵn có 7 chữ số bắt đầu bằng $9999,99$?
\shortans{5}
\loigiai{
\textbf{Đáp số:} $5$\\
$9\,999\,990; 9\,999\,992; 9\,999\,994; 9\,999\,996; 9\,999\,998$ (có 5 số).
}
\end{ex}

\begin{ex} %[4D7V33]
Tìm số bé nhất có 8 chữ số khác nhau có chữ số tận cùng là 8.
\shortans{10234568}
\loigiai{
\textbf{Đáp số:} $10234568$\\
$10\,234\,568$.
}
\end{ex}
